% Pollis Domain Packages: Range-based volatility (RangeVol.jl) -- draft 1 (28 October 2026)
% Compile with: latexmk -pdf RangeVol.tex && latexmk -c
% Shared look for the Pollis New Methods documents (same fonts and colours as the brochure and the workflow handouts).
\documentclass[10pt,a4paper]{article}

\usepackage[a4paper,left=18mm,right=18mm,top=15mm,bottom=17mm,footskip=8mm]{geometry}
\usepackage[T1]{fontenc}
\usepackage[default,semibold]{sourcesanspro}
\usepackage{textcomp}
\usepackage[italic]{mathastext}
\usepackage{amsmath}
\usepackage{microtype}
\usepackage[table]{xcolor}
\usepackage{tikz}
\usepackage[most]{tcolorbox}
\usepackage{fontawesome5}
\usepackage{tabularx,booktabs}
\usepackage{enumitem}
\usepackage{fancyhdr}
\usepackage[hidelinks]{hyperref}

% ---------------------------------------------------------------- colours (as in the brochure)
\definecolor{navy}{HTML}{1C3A5E}
\definecolor{pblue}{HTML}{3B78B5}
\definecolor{amber}{HTML}{F2B33D}
\definecolor{deepamber}{HTML}{D9902A}
\definecolor{pale}{HTML}{EEF3F9}
\definecolor{ink}{HTML}{2A2F36}
\definecolor{logonavy}{HTML}{173D63}
\definecolor{logodark}{HTML}{C98A1E}
\definecolor{logoamber}{HTML}{F2B446}
\definecolor{logocream}{HTML}{FEEFD2}
\definecolor{soft}{HTML}{5E6773}
\definecolor{agree}{HTML}{2F7D4F}
\definecolor{differ}{HTML}{B23B3B}
\color{ink}

\setlength{\parindent}{0pt}
\setlength{\parskip}{4pt}
\setlist{nosep,leftmargin=*,itemsep=2pt}
\renewcommand{\arraystretch}{1.15}

\pagestyle{fancy}
\fancyhf{}
\renewcommand{\headrulewidth}{0pt}
\fancyfoot[L]{\footnotesize\color{soft}\textbf{Pollis}\enspace New Methods\enspace·\enspace\nmtheme}
\fancyfoot[R]{\footnotesize\color{soft}\thepage}
\newcommand{\nmtheme}{}
\newcommand{\nmdate}{11 October 2026}

% The Pollis logo, redrawn so it prints sharply at any size
\newcommand{\pollislogo}[1]{%
  \begin{tikzpicture}[x=#1/512,y=-#1/512,baseline=(current bounding box.center)]
    \fill[logonavy] (0,0) rectangle (512,512);
    \fill[logodark] (36,422) -- (104,256) -- (172,422) -- cycle;
    \fill[logoamber] (104,422) -- (203,180) -- (302,422) -- cycle;
    \fill[logocream] (203,422) -- (339,89) -- (476,422) -- cycle;
  \end{tikzpicture}}

% Top of the first page: logo, name; then the kicker, title and a one-line promise
\newcommand{\nmhead}[3]{%
  \renewcommand{\nmtheme}{#1}%
  \thispagestyle{fancy}%
  \begin{minipage}[c]{13mm}\pollislogo{11mm}\end{minipage}%
  \begin{minipage}[c]{\dimexpr\linewidth-13mm\relax}
    {\color{navy}\Large\bfseries Pollis}\enspace{\color{soft}New Methods}\hfill
    {\color{deepamber}\small\bfseries\MakeUppercase{#1}}
  \end{minipage}\par\vspace{7mm}
  {\color{navy}\fontsize{22}{26}\selectfont\bfseries #2}\par\vspace{4pt}
  {\color{soft}\large #3}\par\vspace{5mm}}

\newcommand{\nmsection}[1]{\par\vspace{5pt}{\color{navy}\large\bfseries #1}\par\vspace{1pt}}
\newcommand{\nmsub}[1]{\par\vspace{2pt}{\color{deepamber}\bfseries #1}\par}
\newcommand{\gt}{\,{\color{deepamber}\bfseries ›}\,}
\newcommand{\menupath}[1]{{\bfseries\color{navy}#1}}
\newcommand{\fn}[1]{\texttt{\small #1}}
% Citations: \nmcite{key}{Breiman (2001)} links to the entry \nmref{key}{...} in the References
\newcommand{\nmcite}[2]{\hyperlink{ref:#1}{\color{pblue}#2}}
\newcommand{\nmref}[2]{\raisebox{1.1em}[0pt][0pt]{\hypertarget{ref:#1}{}}#2\par}
\newcommand{\nmdoi}[1]{\href{https://doi.org/#1}{\color{pblue}doi:#1}}
% Where a notebook is in Pollis: its title in Browse Methods, and its file
\newcommand{\browse}[2]{\menupath{Compose \gt New Methods \gt Browse Methods}: \emph{#1} (\fn{#2.ipynb})}

% Key facts at the top: source, missing in Julia, notebook, result
\newcommand{\factsheet}[4]{%
  \begin{tcolorbox}[enhanced, colback=pale, colframe=pale, boxrule=0pt, sharp corners,
    left=7pt, right=7pt, top=5pt, bottom=5pt, fontupper=\small]
  \begin{tabularx}{\linewidth}{@{}>{\leavevmode\bfseries\color{navy}}p{27mm}>{\leavevmode}X@{}}
  Tested against & #1\\
  Missing in Julia & #2\\
  Notebook & #3\\
  Result & #4
  \end{tabularx}
  \end{tcolorbox}}

% Boxes
\newtcolorbox{pbox}[2][]{enhanced, colback=pale, colframe=amber, boxrule=0pt,
  leftrule=3pt, sharp corners, left=7pt, right=7pt, top=6pt, bottom=6pt,
  fontupper=\small\raggedright, before upper={{\color{navy}\normalsize\bfseries #2}\par\smallskip}, #1}
\newtcolorbox{nbox}[2][]{enhanced, colback=navy, colframe=navy, boxrule=0pt,
  sharp corners, left=8pt, right=8pt, top=7pt, bottom=7pt, coltext=white,
  fontupper=\small\raggedright, before upper={{\color{amber}\normalsize\bfseries #2}\par\smallskip}, #1}

% Value tables: published beside Julia
\newenvironment{valuetable}[1]{%
  \par\small\tabularx{\linewidth}{#1}\toprule}{\bottomrule\endtabularx\par}
\newcommand{\ok}{{\color{agree}\faCheck}}
\newcommand{\no}{{\color{differ}\faTimes}}

% End of the document: credit and licence
\newcommand{\nmcredit}{%
  \par\vfill
  \begin{minipage}[b]{0.78\linewidth}\footnotesize\color{soft}
  Methods written and tested with Claude Code (Anthropic) in Pollis, from the published papers and their data.\par
  \textcopyright\ 2026 Antonio Saragga Seabra\enspace·\enspace Licensed under \href{https://creativecommons.org/licenses/by-sa/4.0/}{CC BY-SA 4.0}\enspace·\enspace Draft, \nmdate
  \end{minipage}\hfill
  \begin{minipage}[b]{0.18\linewidth}\raggedleft\pollislogo{9mm}\par\vspace{2pt}{\footnotesize contact@pollis.org}\end{minipage}}


% A package Pollis needs for a panel of its own, not one that reproduces a paper: its own header, key facts and credit
\renewcommand{\nmdate}{28 October 2026}
\fancyfoot[L]{\footnotesize\color{soft}\textbf{Pollis}\enspace Domain Packages\enspace·\enspace\nmtheme}
\renewcommand{\nmhead}[3]{%
  \renewcommand{\nmtheme}{#1}%
  \thispagestyle{fancy}%
  \begin{minipage}[c]{13mm}\pollislogo{11mm}\end{minipage}%
  \begin{minipage}[c]{\dimexpr\linewidth-13mm\relax}
    {\color{navy}\Large\bfseries Pollis}\enspace{\color{soft}Domain Packages}\hfill
    {\color{deepamber}\small\bfseries\MakeUppercase{#1}}
  \end{minipage}\par\vspace{7mm}
  {\color{navy}\fontsize{22}{26}\selectfont\bfseries #2}\par\vspace{4pt}
  {\color{soft}\large #3}\par\vspace{5mm}}
\renewcommand{\factsheet}[4]{%
  \begin{tcolorbox}[enhanced, colback=pale, colframe=pale, boxrule=0pt, sharp corners,
    left=7pt, right=7pt, top=5pt, bottom=5pt, fontupper=\small]
  \begin{tabularx}{\linewidth}{@{}>{\leavevmode\bfseries\color{navy}}p{27mm}>{\leavevmode}X@{}}
  In Pollis & #1\\
  Missing in Julia & #2\\
  Package & #3\\
  Tests & #4
  \end{tabularx}
  \end{tcolorbox}}
\renewcommand{\nmcredit}{%
  \par\vfill
  \begin{minipage}[b]{0.78\linewidth}\footnotesize\color{soft}
  Package written and tested with Claude Code (Anthropic) in Pollis.\par
  \textcopyright\ 2026 Antonio Saragga Seabra\enspace·\enspace Licensed under \href{https://creativecommons.org/licenses/by-sa/4.0/}{CC BY-SA 4.0}\enspace·\enspace Draft, \nmdate
  \end{minipage}\hfill
  \begin{minipage}[b]{0.18\linewidth}\raggedleft\pollislogo{9mm}\par\vspace{2pt}{\footnotesize contact@pollis.org}\end{minipage}}


\begin{document}
\emergencystretch=3em

\nmhead{Volatility}
  {RangeVol: volatility from the day's high and low}
  {The open, high, low and close of each day say more about volatility than the closes alone: six classical
   estimators that use them, and the log range of Alizadeh, Brandt and Diebold.}

\factsheet
  {The \emph{Measuring Volatility} panel, in Extend > Financial Econometrics, the first panel of the Financial
   Econometrics Toolbox.}
  {the range-based estimators of Parkinson, Garman and Klass, Rogers and Satchell, and Yang and Zhang, with
   close-to-close as the baseline, for daily OHLC prices. R has them in TTR's \fn{volatility()}; Julia had none.}
  {\fn{RangeVol.jl}, from GitHub:\newline \fn{Pkg.add(url = "https://github.com/Trumpingtons/RangeVol.jl")}.}
  {38 tests: 18 against TTR on three datasets (AAPL, the S\&P 500, simulated), 15 on simulated prices, 5 of argument checks. Recomputed today, the six
   estimators agree with TTR 0.24.4 to $2.0\times10^{-14}$ or closer.}

\nmsection{Why a package}
The usual estimate of volatility is the standard deviation of daily returns, close to close. It ignores what happened
during the day: a day that ends where it started may have travelled far in between. \nmcite{parkinson1980}{Parkinson
(1980)} showed that the high and low of the day, in a random walk, estimate the variance about five times as
efficiently as the close; \nmcite{gk1980}{Garman and Klass (1980)} added the open and the close;
\nmcite{rs1991}{Rogers and Satchell (1991)} allowed for a drift; \nmcite{yz2000}{Yang and Zhang (2000)} allowed for
the move overnight, between one close and the next open. Daily OHLC prices are what most data sources give, so these
estimators come at no cost in data. Pollis' panel needed them in Julia, with the same conventions as TTR, so that a
result can be checked in R.

\nmsub{Where they are used}
\begin{itemize}
  \item \textbf{Historical volatility} of stocks, indices, currencies and commodities, for risk measures and option
    pricing, when intraday data are not available.
  \item \textbf{Short windows.} With five times the efficiency, a month of ranges estimates the variance about as
    well as five months of closes.
  \item \textbf{Volatility models.} The log range is a nearly Gaussian measure of volatility for fitting stochastic
    volatility models (\nmcite{abd2002}{Alizadeh, Brandt and Diebold 2002}).
\end{itemize}

\nmsection{The method}
For $n$ days with open $O_t$, high $H_t$, low $L_t$ and close $C_t$, each function returns $\sqrt{\hat\sigma^2 \cdot
k}$, where $\hat\sigma^2$ is the daily variance and $k$ is \fn{trading\_periods}: 1 by default, for daily
volatility; 252 to annualise a stock's, as TTR's \fn{N = 252} does.
\begin{valuetable}{@{}>{\raggedright\arraybackslash}p{52mm}>{\raggedright\arraybackslash}X@{}}
\textbf{Function} & \textbf{Daily variance}\\\midrule
\fn{close\_to\_close(c)} & $\frac1{n-2}\sum_{t=2}^n (r_t - \bar r)^2$, with $r_t = \ln(C_t/C_{t-1})$\\
\fn{parkinson(h, l)} & $\frac{1}{4n\ln 2}\sum_t \ln(H_t/L_t)^2$\\
\fn{garman\_klass(o, h, l, c)} & $\frac1n\sum_t \bigl[\tfrac12\ln(H_t/L_t)^2 - (2\ln2-1)\ln(C_t/O_t)^2\bigr]$\\
\fn{rogers\_satchell(o, h, l, c)} & $\frac1n\sum_t \bigl[\ln(H_t/C_t)\ln(H_t/O_t) + \ln(L_t/C_t)\ln(L_t/O_t)\bigr]$\\
\fn{garman\_klass\_yang\_zhang(o, h, l, c)} & $\frac1{n-1}\sum_{t=2}^n \bigl[\ln(O_t/C_{t-1})^2 + \tfrac12\ln(H_t/L_t)^2 - (2\ln2-1)\ln(C_t/O_t)^2\bigr]$\\
\fn{yang\_zhang(o, h, l, c)} & $\hat\sigma^2_o + \kappa\,\hat\sigma^2_c + (1-\kappa)\,\hat\sigma^2_{RS}$ over $t = 2, \dots, n$, with $\kappa = 0.34/(1.34 + \frac{m+1}{m-1})$, $m = n-1$; needs $n \ge 3$\\
\end{valuetable}
In Yang and Zhang's estimator, $\hat\sigma^2_o$ is the sample variance of the overnight returns $\ln(O_t/C_{t-1})$,
$\hat\sigma^2_c$ that of the open-to-close returns $\ln(C_t/O_t)$, and $\hat\sigma^2_{RS}$ the Rogers--Satchell
variance; $\kappa$ is the weight that minimizes the variance of the estimate. Parkinson's and Garman and Klass's
estimators assume no drift and no move overnight; Rogers and Satchell's allows for a drift; the last two include the
overnight move, and Yang and Zhang's allows for both.

\nmsub{The log range}
\fn{alizadeh\_brandt\_diebold(h, l)} returns a vector: the log range of each day, $\ln(\ln H_t - \ln L_t)$, the
logarithm of the range of the log price (\nmcite{abd2002}{Alizadeh, Brandt and Diebold 2002}). It is close to Gaussian,
with a mean that moves one for one with $\ln\sigma_t$, which makes it an observation of volatility for a state-space
model; a day whose high equals its low gives $-\infty$, to be dropped. Over 100,000 simulated days of 390 steps, the
log range had skewness 0.139 and kurtosis 2.758, with a mean of $0.384 + \ln\sigma$ and a standard deviation of 0.299;
the range itself, \fn{exp.(alizadeh\_brandt\_diebold(h, l))}, had skewness 0.962 and kurtosis 4.176.

\nmsection{A worked example}
\nmsub{The S\&P 500 in 2023}
The \fn{sp500} dataset of PollisDatasets.jl has the daily open, high, low and close of the S\&P 500 since 1985 (Yahoo!
Finance, CC0); 2023 has 250 trading days.
\begin{nbox}{}
\ttfamily\small
using RangeVol, PollisDatasets, Dates\\
d = dataset("sp500"); i = year.(d.date) .== 2023\\
o, h, l, c = d.open[i], d.high[i], d.low[i], d.close[i]\\
yang\_zhang(o, h, l, c)\hfill \# 0.007551, daily\\
yang\_zhang(o, h, l, c; trading\_periods = 252)\hfill \# 0.119866, annualised
\end{nbox}

\begin{valuetable}{@{}lcccccc@{}}
\textbf{Annualised, $k = 252$} & \textbf{close} & \textbf{Parkinson} & \textbf{G--K} & \textbf{R--S} & \textbf{GK--YZ} & \textbf{Y--Z}\\\midrule
S\&P 500 2023, 250 days & 0.130996 & 0.108409 & 0.105824 & 0.103867 & 0.120032 & 0.119866\\
AAPL 2023, 250 days & 0.199154 & 0.178093 & 0.178875 & 0.179683 & 0.213309 & 0.213400\\
\end{valuetable}
For the S\&P 500, the three estimators that see only the trading day give 0.104 to 0.108; those that add the move
overnight, 0.120. For AAPL, from Yahoo! Finance through R's quantmod, the first three give about 0.179 and the last
two 0.213, above close-to-close's 0.199: a single stock moves more between one close and the next open, as earnings
and news released outside trading hours would make it, and the first three miss it.

\nmsub{What the data must be}
The estimators assume that the open, high, low and close come from one session. Before 2006, the open of the
S\&P 500 in Yahoo's data is usually a copy of the previous close, so the move overnight is zero and the open-to-close
return is the whole day's: use the range-based estimators that need the open only from 2008 on. Some sources of
currency prices set the close equal to the open; check before trusting a range-based estimate.

\nmsection{What the tests show}
\nmsub{Against R}
\fn{test/generate\_data.R} downloads AAPL from Yahoo! Finance and the S\&P 500 of 2023 from PollisDatasets.jl, simulates a third series, and saves TTR 0.24.4's
\fn{volatility()} of each with \fn{N = 252}, over the whole series (\fn{n = nrow} for the first four, \fn{n = nrow - 1}
for the two with an overnight term, whose window counts pairs of days). The estimates RangeVol computes today differ
from TTR's by at most $2.0\times10^{-14}$, relative, in each of the 18; running TTR again in R gives the saved values to
about 15 digits. The test suite computes RangeVol's estimates from the saved prices and compares them with TTR's at a
relative tolerance of $10^{-10}$.

\nmsub{By simulation}
\fn{vignette/simulation.jl} draws prices whose log moves as a random walk, with a daily standard deviation of 0.0126,
in 390 steps a day (one a minute) or 23,400 (one a second), with or without a move overnight; 2,000 samples of 250
days, or 400 at seconds. The table gives the mean of each daily variance over the true one, and the efficiency: the
variance of close-to-close's estimate over the estimator's, each divided by its mean.
\begin{valuetable}{@{}lcccccc@{}}
 & \textbf{close} & \textbf{Parkinson} & \textbf{G--K} & \textbf{R--S} & \textbf{GK--YZ} & \textbf{Y--Z}\\\midrule
\multicolumn{7}{@{}l}{\emph{Minutes, no move overnight}}\\
\quad mean / true & 1.003 & 0.936 & 0.910 & 0.910 & 0.910 & 0.923\\
\quad efficiency & 1 & 4.65 & 7.21 & 5.64 & 7.18 & 6.78\\
\multicolumn{7}{@{}l}{\emph{Minutes, a fifth of the variance overnight}}\\
\quad mean / true & 1.001 & 0.749 & 0.728 & 0.728 & 0.928 & 0.938\\
\quad efficiency & 1 & 4.50 & 6.96 & 5.45 & 7.61 & 7.24\\
\multicolumn{7}{@{}l}{\emph{Seconds, no move overnight}}\\
\quad mean / true & 1.001 & 0.994 & 0.991 & 0.991 & 0.991 & 0.992\\
\quad efficiency & 1 & 4.68 & 6.98 & 5.98 & 6.95 & 6.81\\
\end{valuetable}
The efficiencies are close to those the papers derive for continuous prices, about 5.2 for Parkinson's, 7.4 for
Garman and Klass's and 6 for Rogers and Satchell's. A high and low observed once a minute fall short of the true ones,
and the range-based estimators then understate the variance by 6 to 9 per cent; once a second, by less than 1. With
a move overnight, the first three miss it, as in AAPL, and only the last two stay close to the truth.

\nmsub{The rest}
\fn{test/runtests.jl}: 15 tests that each estimator is positive and within 10 per cent of the volatility of 500
simulated days, each drawn in 390 steps with its high and low taken from the path, and that the log range is the
logarithm of the range; and 5 that wrong arguments (vectors of different lengths, too few days) throw.

\nmsection{Still open}
\begin{itemize}
  \item \textbf{Rolling windows.} TTR's \fn{volatility()} returns an estimate for each window of $n$ days; RangeVol
    returns one for the whole series.
  \item \textbf{Bias corrections.} The discrete-sampling bias above has known corrections, and Yang and Zhang's
    $\kappa$ assumes that the drift is small; neither is offered.
\end{itemize}

\nmsection{References}
{\footnotesize
\nmref{abd2002}{Alizadeh, S., Brandt, M.\,W. and Diebold, F.\,X. (2002). Range-based estimation of stochastic volatility models. \emph{Journal of Finance} 57(3), 1047--1091. \nmdoi{10.1111/1540-6261.00454}}
\nmref{gk1980}{Garman, M.\,B. and Klass, M.\,J. (1980). On the estimation of security price volatilities from historical data. \emph{Journal of Business} 53(1), 67--78. \nmdoi{10.1086/296072}}
\nmref{parkinson1980}{Parkinson, M. (1980). The extreme value method for estimating the variance of the rate of return. \emph{Journal of Business} 53(1), 61--65. \nmdoi{10.1086/296071}}
\nmref{rs1991}{Rogers, L.\,C.\,G. and Satchell, S.\,E. (1991). Estimating variance from high, low and closing prices. \emph{Annals of Applied Probability} 1(4), 504--512. \nmdoi{10.1214/aoap/1177005835}}
\nmref{yz2000}{Yang, D. and Zhang, Q. (2000). Drift-independent volatility estimation based on high, low, open, and close prices. \emph{Journal of Business} 73(3), 477--492. \nmdoi{10.1086/209650}}
\nmref{software}{Software: Ulrich, J., TTR, and Ryan, J.\,A. and Ulrich, J.\,M., quantmod (R); PollisDatasets.jl; CSV.jl, Tables.jl, StableRNGs.jl and Statistics (MIT).}
}

\nmcredit
\end{document}
